% TA-PLCA: draft text for the IV 2027 paper (paste into Overleaf; numbers from docs/taplca.md 6.13-6.17).
% Version: TA-PLCA v2 (event anchors, G = 35 us, R = 150 us), campaign of 2026-11-17 (10 seeds x 30 s unless noted).
% Figures: simulations/zonal/t1s_results/figures/fig15_taplca_variation.pdf, fig16_taplca_rtt_ecdf.pdf,
% fig17_taplca_load.pdf. Items marked \todo still need a decision or data.

\subsection{Time-Anchored PLCA}
In standard PLCA, the coordinator emits a BEACON as soon as the last transmit
opportunity (TO) of a cycle ends. The cycle phase therefore drifts with the traffic,
and the time a frame waits for its TO is not known in advance, even when the
release times are synchronized by gPTP. TA-PLCA changes only the coordinator: at the
end of a cycle it holds the BEACON until the next instant of an \emph{anchor list}
$\{a_1, a_2, \dots\}$ (period 10\,ms) on its gPTP-synchronized clock. Each cycle, and
with it each TO, then starts at a known time. Followers run standard PLCA; while they
wait for the BEACON they keep counting TOs, and every hold stays below the 255-TO limit
of the LAN865x ($\approx 816\,\mu$s). If a cycle overruns its anchor, the coordinator
sends the BEACON at once (catch-up).

\paragraph{Event anchors}
A cross-segment frame from segment A to segment B is delivered in the TO~0 of B's
coordinator, so its latency is set by when B's next cycle starts after the frame
arrives. TA-PLCA places one anchor per planned hop. With a bus-timing model measured
in simulation (BEACON 2\,$\mu$s, empty TO 3.2\,$\mu$s, the 64-B frame of follower $i$
ending $2 + 3.2\,i + 62.4\,\mu$s after its anchor, $f_i$), a request/response chain
that starts at $t_0$ gets four anchors:
the client's cycle at $t_0$ (the source releases 15\,$\mu$s earlier);
the far delivery at $d_1 = t_0 + f_i + G$;
the server's reply cycle at $r_1 = d_1 + R$;
and the delivery back at $d_2 = r_1 + f_j + G$.
The gap $G$ absorbs the backbone delay and the clock error, and $R$ is the server's
reply budget. Outside these events, the anchor list is a 25\,$\mu$s grid, so the
coordinator behaves almost like standard PLCA. No grid anchor is allowed within
$G + 30\,\mu$s before an event anchor or during a planned cycle, so no unplanned cycle
can pick up a planned frame early. The lower-priority (TC3) gate on the coordinator
port closes $100 + 5 + G + 30\,\mu$s before each planned anchor (one window cycle plus
margin) and reopens when the planned cycle ends. A sweep over $G$ and $R$ (ideal clocks
and gPTP) gives $G = 35\,\mu$s and $R = 150\,\mu$s as the smallest values that stay
deterministic with gPTP. The planner is generic: periodic chains, one-way flows and
local cycles are listed with their modelled length, and local 1-ms cycles are placed
at the first offset that is free in every millisecond.

\subsection{Evaluation}
All configurations use the dual-bus powertrain scenario with the server processing time measured on
the Raspberry Pi~5 real-time gateway (truncated normal, 2.9\,$\mu$s mean,
$\approx$4.0\,$\mu$s P99), which is included in every round-trip time. We compare
unscheduled releases (random phases); the step-2 schedule with TAS (TC3 closed
400\,$\mu$s per ms), under ideal clocks and under gPTP-synchronized drifting clocks
(U($-$50,\,50)\,ppm, Sync 125\,ms); and TA-PLCA on the same gPTP clocks.
Each configuration is run with 10 seeds $\times$ 30\,s (1\,s warm-up).

\paragraph{Delay variation (H-a)}
Fig.~\ref{fig:taplca-var} shows the delay variation (max$-$min) of the four cross-bus
PCP~7 flows. TA-PLCA reduces it to 0.04--0.12\,$\mu$s, against 102--192\,$\mu$s for the
ideal-clock schedule and 197--503\,$\mu$s for the gPTP schedule. That is three to four
orders of magnitude, and it holds with drifting gPTP clocks. With ideal clocks TA-PLCA
is exactly periodic (0.00\,$\mu$s), so the residual comes from the clock error and the
random processing time. (H-a) holds for all four flows.

\paragraph{Bound (H-b)}
TA-PLCA fixes the round-trip times at 437.41\,$\mu$s (PT01) and 440.41\,$\mu$s (PT02),
processing included (Fig.~\ref{fig:taplca-ecdf}). These are close to the schedules' means
(438--465\,$\mu$s), whose tails reach 848\,$\mu$s, so the bound comes at no cost in mean
latency. The one-way flows PT03/PT04 take 181.41\,$\mu$s, against 227--297\,$\mu$s on
average with the schedules. With the 403\,$\mu$s P99 of the vehicle middleware path as
server time, the gPTP schedule exceeds 1\,ms on every RTT sample (max 1129\,$\mu$s). TA-PLCA
with the same plan ($R = 150\,\mu$s) misses the reply cycle but stays below 1\,ms (max
964\,$\mu$s). With the reply cycle sized for the server ($R = 553\,\mu$s), the round trips are
constant again at 840.41 and 843.41\,$\mu$s (variation $\leq 0.14\,\mu$s), leaving 157\,$\mu$s
below 1\,ms, at the cost of a longer BEACON hold (17\,\% on Z1).

\paragraph{Cost (H-c)}
At nominal load the coordinators hold the BEACON 7\,\% (Z1) and 10\,\% (Z3) of the time,
and the TC3 (window) traffic loses no frames, as with the schedules. The load sweep
(Fig.~\ref{fig:taplca-load}; PT01/PT02 every 10, 5, 2 and 1\,ms) shows that the PCP~7
determinism does not depend on load (worst variation 0.11--0.12\,$\mu$s at every point,
against 503--571\,$\mu$s for the gPTP schedule). The TC3 loss grows only when most
milliseconds carry critical chains: 0, 0, 8 and 70\,\%. On heavier segments (TC3 at twice
the rate, and extra 1-ms PCP~7 local flows; 3 seeds), the gPTP schedule degrades
(mean RTT $+$100--170\,$\mu$s, variation up to 267\,$\mu$s), while TA-PLCA keeps every
cross-bus flow exact. With the doubled TC3 rate it even loses fewer TC3 frames than the
schedule (15 vs.\ 33\,\%). TA-PLCA is therefore a mechanism for segments with idle
capacity, and the claim is bounded by load.

\paragraph{Limitations}
The event anchors leave little slack between a frame's arrival and the far anchor, so
TA-PLCA needs the coordinators' clock error well below $G - 30\,\mu$s: in simulation it
stays exact with gPTP (sub-$\mu$s), but a coordinator offset of 5\,$\mu$s already makes
some hops miss their anchor. Hardware-timestamped PTP ($<$50\,ns on our bench) is far
inside this margin, and software-only synchronization is not. The TC3 gate is part of
the mechanism: without it, the window traffic delays the planned cycles and the bound is
lost. The plan's $G$ and $R$ were chosen by a sweep, not by a solver. Compatibility with
standard followers that wait for a deferred BEACON must still be confirmed against the
published Clause~148 text and on a LAN865x.

% \begin{figure}[t]\centering
% \includegraphics[width=\columnwidth]{fig15_taplca_variation.pdf}
% \caption{Delay variation (max$-$min) of the cross-bus PCP~7 flows, nominal load, 10 seeds $\times$ 30\,s (log scale).}
% \label{fig:taplca-var}\end{figure}
% \begin{figure}[t]\centering
% \includegraphics[width=\columnwidth]{fig17_taplca_load.pdf}
% \caption{Load sweep: worst delay variation of the cross-bus PCP~7 flows (left, log scale) and TC3 frame loss (right), 10 seeds $\times$ 30\,s per point.}
% \label{fig:taplca-load}\end{figure}
% \begin{figure}[t]\centering
% \includegraphics[width=\columnwidth]{fig16_taplca_rtt_ecdf.pdf}
% \caption{ECDF of the request/response round-trip times, processing included.}
% \label{fig:taplca-ecdf}\end{figure}
